% SEGMENT OLP-0582-S01
\documentclass[../../../include/open-logic-section]{subfiles}

\begin{document}

\olfileid{sth}{cardinals}{cardsasords}
\olsection{வரிசையெண்களாகக் கண அளவெண்கள்}

உண்மையில், கண அளவெண்களைப் பற்றிய நமது கோட்பாடு வரிசையெண்களைப்
பற்றிய கோட்பாட்டை (எந்தத் தயக்கமும் இன்றி) பயன்படுத்தும். அதாவது,
குறிப்பிட்ட சில வரிசையெண்களையே கண அளவெண்களாக வரையறுப்போம்.
குறிப்பாக, பின்வரும் வரையறையை அளிக்கிறோம்:

\begin{defn}\ollabel{defcardinalasordinal}
$A$ ஐ நன்கு வரிசைப்படுத்த முடிந்தால், $\cardeq{A}{\gamma}$ ஆகுமாறு
உள்ள மீச்சிறிய வரிசையெண் $\gamma$ தான் $\card{A}$. எந்த வரிசையெண்
$\gamma$ க்கும், $\gamma = \card{\gamma}$ எனும்போதும் அப்போதுதான்
$\gamma$ ஒரு \emph{கண அளவெண்} என்று சொல்கிறோம்.
\end{defn}

``$A$ ஐ நன்கு வரிசைப்படுத்த முடியும்'' என்ற சொற்றொடரை இப்போது
பயன்படுத்தினோம். கணிதத்தில் பெரும்பாலும் இருப்பதைப் போல, இங்குள்ள
சாத்தியத்தைக் குறிக்கும் சொற்கள் ஓர் இருப்புக் கூற்றை முறைசாராமல்
குறிக்கின்றன: ``$A$ ஐ நன்கு வரிசைப்படுத்த முடியும்'' என்பதன் பொருள்,
``$A$ ஐ நன்கு வரிசைப்படுத்தும் ஒரு தொடர்பு உள்ளது'' என்பதே.

ஆனால் \olref{defcardinalasordinal} இல் ஒரு சிக்கல் உள்ளது.
\emph{ஒவ்வொரு} கணத்திற்கும் ஓர் அளவு இருக்க வேண்டும்; அதாவது,
ஒவ்வொரு $A$ க்கும் $\card{A}$ இருக்க வேண்டும் என்று விரும்புகிறோம்.
எனினும், நாம் இப்போது தந்த வரையறை ஒரு நிபந்தனையுடன் தொடங்குகிறது:
``$A$ ஐ நன்கு வரிசைப்படுத்த \emph{முடிந்தால்}\ldots''.
நன்கு வரிசைப்படுத்த முடியாத ஒரு கணம் $A$ இருந்தால், அந்த வரையறை
$\card{A}$ என்ற பொருளை வரையறுக்காது.

ஆகவே, \olref{defcardinalasordinal} ஐப் பயன்படுத்த, ஒவ்வொரு கணத்தையும்
நன்கு வரிசைப்படுத்த முடியும் என்ற உறுதி தேவை. ஆனால் $\ZF$ இல்
அந்த உறுதி கிடைக்கவில்லை. எனவே, \olref{defcardinalasordinal} ஐப்
பயன்படுத்த விரும்பினால், பின்வருவது போன்ற ஒரு புதிய அடிகோளைச்
சேர்ப்பதைத் தவிர வேறு வழியில்லை:
\begin{axiom}[நன்கு வரிசைப்படுத்தல்]
ஒவ்வொரு கணத்தையும் நன்கு வரிசைப்படுத்த முடியும்.
\end{axiom}
இந்த நன்கு வரிசைப்படுத்தல் அடிகோள் ஏற்கத்தக்கதா என்பதை
\olref[choice][]{chap} இல் ஆராய்வோம். ஆனால் இனி அதை நாம்
பயன்படுத்திக் கொள்வோம்.

% SEGMENT OLP-0582-S02
அதைப் பயன்படுத்தினால், \olref{defcardinalasordinal} இல் வரையறுக்கப்பட்ட
கண அளவெண்கள் இருக்கின்றன என்றும் சரியாகச் செயல்படுகின்றன என்றும்
நிறுவுவது எளிது:

\begin{lem}\ollabel{lem:CardinalsExist}
ஒவ்வொரு கணம் $A$ க்கும்:
\begin{enumerate}
	\item\ollabel{cardaexists} $\card{A}$ இருக்கிறது; அது தனித்துவமானது;
	\item\ollabel{cardaapprox}  $\cardeq{\card{A}}{A}$;
	\item\ollabel{cardaidem}  $\card{A}$ ஒரு கண அளவெண்; அதாவது,
	$\card{A} = \card{\card{A}}$;
\end{enumerate}
\end{lem}

\begin{proof}
$A$ ஐ நிலையாகக் கொள்வோம். நன்கு வரிசைப்படுத்தல் அடிகோளின்படி,
$\tuple{A, R}$ என்ற நன்கு வரிசைப்படுத்தல் உள்ளது.
\olref[ordinals][ordtype]{thmOrdinalRepresentation} இன்படி,
$\tuple{A, R}$ ஆனது $\beta$ என்ற தனித்துவமான வரிசையெண்ணுடன்
வரிசை ஓரின அமைப்புடையது. எனவே $\cardeq{A}{\beta}$.
முடிவிலிக்கடந்த தொகுத்தறிதலின்படி, $\cardeq{A}{\gamma}$ ஆகுமாறு
உள்ள தனித்துவமான மீச்சிறிய வரிசையெண் $\gamma$ உள்ளது. ஆகவே
$\card{A} = \gamma$; இதனால் \olref{cardaexists},
\olref{cardaapprox} ஆகியவை நிறுவப்பட்டன.
\olref{cardaidem} ஐ நிறுவ, $\delta \in \gamma$ எனில்,
$\gamma$ ஐத் தேர்ந்தெடுத்த விதத்தால் $\cardless{\delta}{A}$
என்பதைக் கவனிக்கவும். எண்ணொத்த தன்மை ஒரு சமானத் தொடர்பு
(\olref[sfr][siz][equ]{equinumerosityisequi}) என்பதால்,
$\cardless{\delta}{\gamma}$ என்றும் பெறுகிறோம். எனவே
$\gamma = \card{\gamma}$.
% மறுதலித்து நிறுவ, $\delta \in \gamma$ மற்றும் $\cardeq{\gamma}{\delta}$ ஆகுமாறு ஒரு வரிசையெண் $\delta$ உள்ளது எனக் கொள்வோம். அப்போது $\cardeq{\cardeq{A}{\gamma}}{\delta}$; ஆகவே \olref[sfr][set][equ]{equinumerosityisequi} இன்படி $\cardeq{A}{\delta}$, இது $\gamma$ இன் மீச்சிறிய தேர்வுக்கு முரணாகும்.
\end{proof}

% SEGMENT OLP-0582-S03
அடுத்த முடிவு காண்டோரின் கொள்கையையும் மேலும் பலவற்றையும் உறுதி செய்கிறது.
(கண அளவெண்கள் தமது வரிசையை வரிசையெண்களிலிருந்து பெறுகின்றன என்பதைக்
கவனிக்கவும்; அதாவது $\cardfont{a} < \cardfont{b}$ எனும்போதும்
அப்போதுதான் $\cardfont{a} \in \cardfont{b}$. இதைக் கூறும்போது,
கண அளவெண்கள் என்று அறிந்துள்ள பொருள்களுக்கு ஃபிராக்டூர் எழுத்துகளைப்
பயன்படுத்துவோம். இது பொதுவான மரபு. கண அளவெண்கள் வரிசையெண்களாகவும்
இருப்பதால், கிரேக்க எழுத்துமாலையின் நடுப்பகுதி எழுத்துகளைத் தேர்வது
மற்றொரு பொதுவான முறை; எடுத்துக்காட்டாக, $\kappa, \lambda$.)
\begin{lem}\ollabel{lem:CardinalsBehaveRight}
எந்தக் கணங்கள் $A$, $B$ க்கும்:
\begin{align*}
	\cardeq{A}{B} &\text{ iff } \card{A} = \card{B}\\
	\cardle{A}{B} &\text{ iff } \card{A} \leq \card{B}\\
	\cardless{A}{B}&\text{ iff } \card{A} < \card{B}
\end{align*}
\end{lem}

\begin{proof}
இரண்டாவது கூற்றின் இடமிருந்து வலம் செல்லும் திசையை நிறுவுவோம்;
மற்ற நிலைகளும் இதைப் போன்றவை, அவை பயிற்சிக்கு விடப்படுகின்றன.
பின்வரும் படத்தைக் கருதுங்கள்:
\begin{center}
	\begin{tikzpicture}
	\node (nodea) {$A$};
	\node[right = 6em of nodea] (nodeb) {$B$};
	\node[below = 2em of nodea] (nodecarda) {$\card{A}$};
	\node[below = 2em of nodeb] (nodecardb) {$\card{B}$};
	\draw[->] (nodea)--(nodeb);
	\draw[<->] (nodea)--(nodecarda);
	\draw[<->] (nodeb)--(nodecardb);
	\draw[->, dashed] (nodecarda)--(nodecardb);
\end{tikzpicture}
\end{center}
இரு முனைகளிலும் அம்புள்ள கோடுகள் !!{bijection}s ஐக் குறிக்கின்றன;
\olref{lem:CardinalsExist} அவை இருப்பதை உறுதி செய்கிறது.
$\cardle{A}{B}$ என்று கருதுவதால், $A\to B$ என்ற
!!a{injection} உள்ளது. இப்போது $\card{A}$ இலிருந்து $A$ க்கும்,
அங்கிருந்து $B$ க்கும், பிறகு $\card{B}$ க்கும் அம்புகளைத் தொடர்ந்து,
$\card{A} \to \card{B}$ என்ற !!a{injection} ஐப் பெறுகிறோம்
(இது கோடிட்ட அம்பு).
\end{proof}\noindent
\olref{lem:CardinalsBehaveRight} ஐப் பயன்படுத்தி
ஷ்ரோடர்--பெர்ன்ஸ்டைன் தேற்றத்தை மீண்டும் நிறுவலாம்.
$\cardle{A}{B}$ மற்றும் $\cardle{B}{A}$ எனில் $\cardeq{A}{B}$
என்பதே அக்கூற்று. இதை \olref[sfr][siz][sb]{thm:schroder-bernstein}
என்று கூறினோம்; ஆனால் முதலில் \olref[sfr][infinite][card-sb]{sec} இல்
சற்று முயற்சியுடன் நிறுவினோம். இப்போது பாருங்கள்:

\begin{proof}[ஷ்ரோடர்--பெர்ன்ஸ்டைனின் மறு நிறுவல்]
$\cardle{A}{B}$ மற்றும் $\cardle{B}{A}$ எனில்,
\olref{lem:CardinalsBehaveRight} இன்படி $\card{A} \leq \card{B}$
மற்றும் $\card{B} \leq \card{A}$. ஆகவே, முக்கூற்றுப் பண்பினாலும்
\olref{lem:CardinalsBehaveRight} இனாலும் $\card{A} = \card{B}$,
$\cardeq{A}{B}$ ஆகியவை கிடைக்கின்றன.
\end{proof}
\noindent
இந்த நிறுவல் மிக எளிதாக இருந்தாலும், இதன் பின்னணியில்
\olref[ordinals][ordtype]{thmOrdinalRepresentation} ஐப் பெற
மாற்றீட்டு அடிகோளும், \olref{lem:CardinalsBehaveRight} ஐ உறுதிசெய்ய
நன்கு வரிசைப்படுத்தல் அடிகோளும் பயன்படுத்தப்படுகின்றன. மாறாக,
\olref[sfr][infinite][card-sb]{sec} இன் நிறுவல் அதிகம் தன்னிறைவு
கொண்டது; அதனை $\Zminus$ இலும் செய்ய முடியும்.

\end{document}
